\documentclass[11pt,aspectratio=169]{beamer}

\usetheme{Madrid}
\usecolortheme{seahorse}

\usepackage[utf8]{inputenc}
\usepackage{amsmath,amssymb}
\usepackage{xcolor}
\usepackage{booktabs}
\usepackage{tikz}
\usetikzlibrary{arrows.meta,positioning,fit,backgrounds}

\setbeamertemplate{navigation symbols}{}
\setbeamertemplate{footline}[frame number]

\title{Beyond the Syllabus}
\subtitle{Randomness, Refined: Karger--Stein, Derandomization, and the Limits of
Coin-Flipping\\
\small Week 15 Supplement -- TOAE209/TOAH209: Design and Analysis of Algorithms}
\author{Foreign Trade University}
\date{}

\begin{document}

% ---------------------------------------------------------------
\begin{frame}
  \titlepage
\end{frame}
% ---------------------------------------------------------------

\begin{frame}{Outline}
  \tableofcontents
\end{frame}

% =================================================================
\section{Motivation}
% =================================================================

\begin{frame}{Why look beyond the syllabus?}
  \small
  \begin{itemize}
    \item This week you met five randomized algorithms and the Las Vegas/Monte Carlo
    distinction: QuickSort ($O(n\log n)$ expected), Quickselect ($O(n)$ expected),
    Karger's contraction algorithm for global min-cut, reservoir sampling, and
    universal hashing / Miller--Rabin -- all analyzed with indicator variables,
    linearity of expectation, and the union bound.
    \item Three questions a curious student should ask next:
    \begin{enumerate}
      \item Karger's algorithm, as amplified in lecture, reruns from scratch
      $\Theta(n^2\log n)$ times. Can the \emph{same} probability argument you proved be
      reused more cleverly instead of just repeated?
      \item Is randomness a genuinely \textbf{separate} computational resource from
      determinism -- or can every efficient randomized algorithm, in principle, be
      \textbf{derandomized} at only polynomial cost?
      \item For Quickselect specifically, has that derandomization question already
      been \emph{settled}?
    \end{enumerate}
  \end{itemize}
  \begin{alertblock}{How this supplement is different from a survey}
    Every section below walks through the \emph{actual mechanics} of a result --
    reusing a proof technique you already hand-traced this week -- not just a headline
    number. Every specific bound is sourced to a specific paper.
  \end{alertblock}
\end{frame}

% =================================================================
\section{Sharpening This Week's Own Algorithm: Karger--Stein}
% =================================================================

\begin{frame}{Recap: the cost of amplifying Karger's algorithm}
  \begin{block}{What you proved this week}
    A single contraction run keeps a fixed min-cut with probability
    $\ge 1/\binom{n}{2}$ (Theorem, lecture notes). Amplifying by
    $T = \binom{n}{2}\ln n = \Theta(n^2\log n)$ \textbf{independent, from-scratch}
    reruns drives the failure probability down to $\le 1/n$.
  \end{block}
  \begin{itemize}
    \item Each full run costs $\Theta(m)$ (one pass of $n-2$ edge contractions on a
    graph with $m$ edges), so the amplified scheme costs
    $\Theta(mn^2\log n)$ total.
    \item That is polynomial -- but it treats every rerun as throwing away \emph{all}
    prior work and starting over from $n$ vertices. Is that wasteful?
  \end{itemize}
\end{frame}

\begin{frame}{The key idea (Karger \& Stein, JACM 1996): recurse, don't restart}
  \begin{block}{Karger \& Stein, ``A New Approach to the Minimum Cut Problem,''
  \emph{J.\ ACM} 43(4):601--640, 1996}
    Reuse the \emph{exact same} telescoping survival-probability argument from this
    week's lecture -- but stop it early, and recurse.
  \end{block}
  \begin{itemize}
    \item The lecture's proof shows: contracting from $n$ down to $n'$ vertices keeps a
    fixed min-cut with probability $\ge \dfrac{n'(n'-1)}{n(n-1)}$ (same telescoping
    product, just not carried all the way to $2$).
    \item Set $n' = n/\sqrt2$: that ratio is $\approx 1/2$. So contracting \emph{halfway}
    (in this $\sqrt2$ sense) survives with probability $\ge 1/2$ -- much better odds
    than surviving all the way to $2$ vertices.
    \item \textbf{Recursive algorithm:} contract $n \to n/\sqrt2$ once, then recurse
    \textbf{twice}, independently, on that smaller graph, and keep the better of the
    two results. Base case: a small constant-size graph.
  \end{itemize}
\end{frame}

\begin{frame}{Solving the two recurrences}
  \small
  \begin{block}{Time}
    $T(n) = 2\,T(n/\sqrt2) + O(n^2)$ (the $O(n^2)$ is the cost of one contraction phase
    down to $n/\sqrt2$ vertices). Since $\log_{\sqrt2} 2 = 2$, the branching exactly
    balances the shrink rate at every level, giving
    \[
      T(n) = O(n^2 \log n).
    \]
  \end{block}
  \begin{block}{Success probability}
    $P(n) = 1 - \bigl(1 - \tfrac12 P(n/\sqrt2)\bigr)^2$ (fails only if \emph{both}
    recursive branches fail). Solving this recurrence gives
    \[
      P(n) = \Theta(1/\log n).
    \]
  \end{block}
  \begin{exampleblock}{The payoff}
    Plain Karger: one run succeeds with probability $\Theta(1/n^2)$. One
    Karger--Stein call succeeds with probability $\Theta(1/\log n)$ -- exponentially
    better per run, for only an extra $\log n$ factor in the per-run cost.
  \end{exampleblock}
\end{frame}

\begin{frame}{Putting it together: amplifying the better algorithm}
  \small
  \begin{itemize}
    \item Exactly the union-bound amplification argument from this week's Corollary
    (Karger amplification) and Theorem (contention resolution) applies again: to drive
    failure probability down to $\le 1/n$ with a per-run success probability of
    $\Theta(1/\log n)$, you need only $O(\log n)$ independent reruns instead of
    $O(n^2\log n)$.
    \item Total cost: $O(\log n)$ reruns $\times\ O(n^2\log n)$ per run
    $= O(n^2\log^3 n)$ -- exactly the bound Karger and Stein proved.
  \end{itemize}
  \begin{center}
  \begin{tabular}{@{}l c p{4.3cm}@{}}
    \toprule
    Scheme & Total time & Idea \\
    \midrule
    Plain Karger, amplified (lecture) & $\Theta(mn^2\log n)$ & $\Theta(n^2\log n)$
    independent full reruns \\
    \textbf{Karger--Stein (1996)} & $O(n^2\log^3 n)$ & recurse on $2$ halved copies,
    reuse the same proof \\
    \bottomrule
  \end{tabular}
  \end{center}
  For sparse graphs ($m \approx n$) this removes an entire factor of $n$ -- no new
  probability idea, just reusing the lecture's own telescoping argument \emph{locally}
  instead of only globally.
\end{frame}

\begin{frame}{Is Karger--Stein itself optimal?}
  \begin{block}{Gupta, Lee, and Li, ``The Karger--Stein Algorithm Is Optimal for
  $k$-Cut,'' STOC 2020}
    The same recursive-contraction strategy, applied to the more general problem of
    splitting a graph into $k$ pieces (not just $2$), outputs any fixed minimum $k$-cut
    with probability that matches the best possible bound, up to lower-order terms.
  \end{block}
  \begin{itemize}
    \item Same theme as Week 3's ``is Union-Find provably optimal'' story: a simple,
    20-line randomized idea from the 1990s turns out to be not just good, but
    (nearly) the best \emph{any} algorithm of its kind can do -- confirmed 24 years
    later, and for a strictly harder generalization of the problem.
  \end{itemize}
\end{frame}

% =================================================================
\section{Can Randomness Always Be Removed?}
% =================================================================

\begin{frame}{The question: is BPP = P?}
  \small
  \begin{itemize}
    \item \textbf{BPP} is the class of problems solvable by a Monte Carlo algorithm
    running in polynomial time with bounded error (exactly this week's Monte Carlo
    definition, formalized as a complexity class).
    \item Open since the 1980s: is \textbf{every} such algorithm derandomizable --
    replaceable by a deterministic polynomial-time algorithm? I.e., is $BPP = P$?
    \item The classical attack, \textbf{hardness-vs-randomness} (Nisan \& Wigderson,
    \emph{J.\ Comput.\ Syst.\ Sci.} 1994): \emph{if} some function needs large circuits
    to compute in the worst case, \emph{then} that same hardness can be converted into a
    pseudorandom generator (PRG) -- a short deterministic seed that ``looks random''
    enough to fool any small circuit, replacing true coin flips at only polynomial cost.
  \end{itemize}
  \begin{alertblock}{Why this matters here}
    If this program succeeds completely, every algorithm in this week's lecture --
    QuickSort, Quickselect, Karger's contraction -- would have a deterministic
    polynomial-time counterpart with (asymptotically) the same running time. That is a
    long-standing \emph{research program}, not a settled fact.
  \end{alertblock}
\end{frame}

\begin{frame}{A landmark tightening (2020/2022)}
  \begin{block}{Doron, Moshkovitz, Oh, and Zuckerman, ``Nearly Optimal Pseudorandomness
  from Hardness,'' FOCS 2020 / \emph{J.\ ACM} 69(6), Article 43, 2022 --
  \emph{DOI:} \texttt{10.1145/3555307}}
    Construct pseudorandom generators whose parameters (how hard the seed function must
    be vs.\ how large a circuit the generator fools) are \textbf{nearly optimal} for the
    hardness-vs-randomness paradigm -- closing most of a 26-year-old gap between what
    Nisan--Wigderson-style constructions could prove and what the paradigm should, in
    principle, be able to deliver.
  \end{block}
  \begin{itemize}
    \item This does not resolve $BPP \stackrel{?}{=} P$ (it is still conditional on an
    unproven worst-case hardness assumption) -- but it means \emph{if} that assumption
    holds, the derandomization it buys you is essentially as good as this framework can
    ever give.
  \end{itemize}
\end{frame}

\begin{frame}{Still being tightened -- this month}
  \begin{block}{Doron, Moshkovitz, Oh, and Zuckerman, ECCC Technical Report TR26-082,
  May 2026}
    The same four authors continue optimizing exactly this tradeoff: how much
    worst-case circuit hardness you must assume vs.\ how large a circuit the resulting
    generator can fool, pushing the two parameters closer together than the 2022 result.
  \end{block}
  \begin{alertblock}{An open program, moving in real time}
    As of this month (August 2026), the parameters of a 30-year-old paradigm --
    hardness-vs-randomness -- are still being sharpened. This is not settled
    textbook history; it is active research you could, in principle, read the newest
    preprint of.
  \end{alertblock}
\end{frame}

\begin{frame}{One case where derandomization already won -- in 1973}
  \small
  \begin{itemize}
    \item The \emph{general} $BPP \stackrel{?}{=} P$ question is open, but individual
    problems can be -- and have been -- derandomized completely.
    \item This week's own \textbf{Quickselect} (expected $O(n)$, worst-case
    $\Theta(n^2)$) has a fully deterministic worst-case-$O(n)$ counterpart, found
    \emph{47 years before} the 2020 pseudorandomness result above:
  \end{itemize}
  \begin{block}{Blum, Floyd, Pratt, Rivest, and Tarjan, ``Time Bounds for Selection,''
  \emph{J.\ Comput.\ Syst.\ Sci.} 7:448--461, 1973}
    The \textbf{median-of-medians} algorithm: split into groups of $5$, recursively find
    the median of each group's median, and use \emph{that} as a provably good pivot --
    no coin flips anywhere, worst-case $O(n)$, guaranteed on every input.
  \end{block}
  \begin{exampleblock}{The lesson}
    ``Open in general'' does not mean ``open for the problem you care about.'' Median-of-medians
    is asymptotically optimal but has large constants, which is exactly why randomized
    Quickselect -- simpler, faster in practice -- is what you actually traced this week.
  \end{exampleblock}
\end{frame}

% =================================================================
\section{Randomness Is Still Winning, Right Now}
% =================================================================

\begin{frame}{The current fastest max-flow algorithm is randomized}
  \begin{block}{Chen, Kyng, Liu, Peng, Probst Gutenberg, and Sachdeva, ``Maximum Flow
  and Minimum-Cost Flow in Almost-Linear Time,'' FOCS 2022 (arXiv:2203.00671)}
    Computes exact maximum flow and minimum-cost flow on graphs with $m$ edges and
    polynomially bounded integer data in $m^{1+o(1)}$ time -- a landmark speedup over
    the polynomial (but far slower) flow algorithms in a typical course.
  \end{block}
  \begin{itemize}
    \item The algorithm's numerical core relies on \textbf{randomized} graph
    sparsification and randomized linear-algebra solvers -- it is not just historically
    related to this week's ideas, it is \emph{built from} the same coin-flipping toolkit.
    \item As of August 2026, no deterministic algorithm matches this speed. Randomness
    is not a classroom curiosity here; it is the current frontier.
  \end{itemize}
\end{frame}

% =================================================================
\section{Summary}
% =================================================================

\begin{frame}{Summary}
  \footnotesize
  \begin{enumerate}
    \item \textbf{Even the algorithm you hand-traced had headroom left:}
    Karger--Stein (1996) reuses this week's own telescoping-probability proof
    \emph{locally} instead of only globally, cutting an amplified $\Theta(mn^2\log n)$
    scheme down to $O(n^2\log^3 n)$ -- and that recursive strategy was proven
    (near-)optimal for the general $k$-cut problem in 2020.
    \item \textbf{Is randomness a separate resource from determinism?} The
    hardness-vs-randomness program (Nisan--Wigderson, 1994) is still being tightened --
    most recently this month (Doron, Moshkovitz, Oh, Zuckerman, May 2026) -- with no
    final answer yet.
    \item \textbf{But some cases are already closed:} Quickselect's own worst case was
    fully derandomized in 1973 (median-of-medians), decades before the general question
    was even well-posed.
    \item \textbf{Randomization is still winning at the frontier:} the fastest known
    max-flow algorithm (2022) is randomized at its numerical core, with nothing
    deterministic matching it as of today.
  \end{enumerate}
  \begin{alertblock}{Looking back over the course}
    The lecture's own closing section traces six paradigms across fifteen weeks. This
    supplement's point: several of them -- greedy's exchange argument, flow's
    augmenting paths, this week's contraction algorithm -- were sharpened within the
    last five years by exactly the move Karger--Stein made: take a proof you trust, and
    reuse it more cleverly.
  \end{alertblock}
\end{frame}

\end{document}
